% !TeX root = ../paper.tex

\section{Methodology}
\label{sec:method}

\subsection{Study Design}
The study uses a quantitative, experimental design with paired measurements. For
each task, a no-feedback baseline is compared against the same task after up to
three static-analysis feedback iterations. Pairing each baseline with its
refined counterpart controls for task-level difficulty.

\subsection{Materials and Tools}
The target language is Python. Generation uses
\textbf{Qwen2.5-Coder-7B-Instruct}~\citep{qwen25coder}, an open-weight code model,
served locally by Ollama~\citep{ollama} under the \texttt{qwen2.5-coder:7b} tag
(Q4\_K\_M quantization), with temperature~0 and a fixed seed for reproducibility. This model was chosen
deliberately: a mid-tier, open-weight code model is strong enough to give a
meaningful correctness baseline yet still leaves structural headroom for the
feedback loop to act on. Local hosting was chosen so
that a pinned model, quantization, and seed make generation reproducible without
provider drift and at no API cost. Each metric maps to one open-source tool, as
summarized in \cref{tab:tooling}.

\begin{table*}[t]
  \centering
  \caption{Metrics and the tools that produce them. Tools:
    EvalPlus~\citep{evalplus_tool}, \texttt{radon}~\citep{radon_docs},
    \texttt{complexipy}~\citep{complexipy_tool}, \texttt{pylint}~\citep{pylint_tool}.}
  \label{tab:tooling}
  \begin{tabular}{lll}
    \toprule
    \textbf{Metric} & \textbf{Role} & \textbf{Tool} \\
    \midrule
    Functional correctness            & Primary (H3) & EvalPlus (MBPP+) \\
    Maintainability Index             & Primary (H1) & \texttt{radon} \\
    Cyclomatic complexity             & Primary (H2) & \texttt{radon} \\
    Cognitive complexity              & Primary (H2) & \texttt{complexipy} \\
    Halstead volume/effort            & Supplementary & \texttt{radon} \\
    Lines of code (LOC)               & Supplementary / control & \texttt{radon} \\
    Linter counts (R/W/C)             & Feedback source & \texttt{pylint} \\
    \bottomrule
  \end{tabular}
\end{table*}

\noindent Pylint is positioned strictly as a feedback mechanism, not as a primary
scientific measure. Because the loop is shown its own Pylint findings, any later
movement in those counts partly restates the input it was given, so they are
reported as descriptive context only (\cref{sec:supp}). Results are stored as
JSONL, one record per task and iteration, alongside the generated code.

\subsection{Feedback Pipeline Architecture}
% The supervisor noted the loop was previously under-explained. Each stage below
% is described in one short step; the figure mirrors these stages 1:1.
The pipeline runs one task at a time through the following stages
(\cref{fig:pipeline}):
\begin{enumerate}
  \item \textbf{Generate.} The LLM receives the natural-language task prompt and
    produces an initial Python solution.
  \item \textbf{Check correctness.} The solution is executed against the MBPP+
    inputs and its outputs are compared to the reference (canonical) solution used
    as an oracle; the pass/fail outcome is recorded.
  \item \textbf{Analyze.} The static-analysis tools parse the solution and emit
    the metric values and specific violations.
  \item \textbf{Format feedback.} The measured values and violations are rendered
    into a structured feedback prompt, reproduced in full in \cref{fig:prompt}.
  \item \textbf{Refine.} The LLM receives its previous code plus the feedback and
    returns a refactored solution.
\end{enumerate}
All of this happens in the prompt. The model's weights are never updated: there is
no fine-tuning and no gradient step, and the static-analysis tools are run by the
pipeline rather than called by the model. ``Feedback'' throughout this thesis means
tool output formatted into the next prompt, together with the model's previous
code.

Because that prompt is the study's intervention, it is given verbatim in
\cref{fig:prompt} rather than paraphrased. The metric values shown there are
illustrative; everything else is the exact template, and the three configurations
of \cref{sec:configs} are defined by small, stated edits to it.

\begin{figure}[t]
\begin{lstlisting}[basicstyle=\ttfamily\scriptsize,numbers=none,language=,
                   frame=single,breaklines=true,columns=fullflexible]
Your previous Python solution was:
```python
{previous solution}
```

A static-analysis tool reported the following on
that solution:
- Maintainability Index: 78.2 (0-100, higher is
  better)
- Cyclomatic complexity (total): 7 (independent
  branches; lower is simpler)
- Cognitive complexity (total): 12 (penalises
  deep/nested control flow; lower is easier to
  follow)
- Pylint findings:
    line 4: too-many-branches -- Too many
    branches (14/12)

Refactor only the internal logic of the function
to reduce its cyclomatic and cognitive
complexity: flatten nested loops and
conditionals, remove redundant branches, and use
clear idiomatic Python (comprehensions,
built-ins, early returns) where it genuinely
simplifies the control flow. Keep the EXACT same
function signature (same name AND parameters)
and the same behaviour. Keep any existing
docstring and comments -- do not delete them.
Original task:
{task description}

Respond with ONLY a Python code block.
\end{lstlisting}
\caption{The refactor prompt, verbatim, for the comment-preserving configuration.
  Metric values are illustrative. The comment-stripping configuration omits the
  final sentence of the instruction paragraph and is otherwise identical; the
  correctness-guarded configuration uses this template unchanged and prepends a
  rejection notice after a refused refactor; the code-blind variant drops the
  previous solution and the line-anchored Pylint findings, keeping only the scalar
  values.}
\label{fig:prompt}
\end{figure}

Stages 2--5 repeat for \emph{up to three} iterations. The cap is a fixed budget set
in advance: it gives the loop room to make successive improvements while keeping
every run bounded, mutually comparable, and cheap to reproduce. The loop does not stop early on convergence; it always runs the full three
iterations, so the number of refinements is identical across tasks and configurations. Correctness is
re-checked on every iteration (Stage~2), so any refactor that breaks a test is
detected immediately; how the loop then responds---carry the new version forward
regardless, or reject it and roll back---is the one axis on which the
configurations differ (\cref{sec:configs}).

\begin{figure*}[t]
  \centering
  \begin{tikzpicture}[
      node distance=10mm and 14mm,
      box/.style={rectangle,draw,rounded corners,align=center,
                  minimum height=9mm,minimum width=26mm,font=\small},
      arr/.style={-{Stealth[length=2.5mm]}}]
    \node[box] (gen)   {LLM generates\\Python code};
    \node[box,right=of gen] (analyze) {Static analysis\\(radon, pylint)};
    \node[box,right=of analyze] (feedback) {Format report\\as feedback prompt};
    \node[box,below=of feedback] (refine) {LLM refines\\code};
    \node[box,below=of gen] (test) {MBPP+ oracle\\correctness check};
    \draw[arr] (gen) -- (analyze);
    \draw[arr] (analyze) -- (feedback);
    \draw[arr] (feedback) -- (refine);
    \draw[arr] (refine) -- node[midway,above=0.5mm,font=\scriptsize]{up to 3$\times$} (test);
    \draw[arr] (test) -- (gen);
  \end{tikzpicture}
  \caption{Iterative static-analysis feedback pipeline.}
  \label{fig:pipeline}
\end{figure*}

\subsection{Task Set and Correctness Oracle}
\label{sec:taskset}
Tasks are drawn from the sanitized MBPP problems~\citep{austin2021mbpp} and are
evaluated with MBPP+. EvalPlus~\citep{liu2023evalplus} is a framework that augments
a benchmark with a substantially larger, automatically generated set of test
inputs; the paper applies it to HumanEval and names MBPP as future work. MBPP+ is
the result of that later application and is distributed with the EvalPlus tool,
version~0.3.1~\citep{evalplus_tool}, which is the version used here. The augmented
suite makes the correctness check considerably more rigorous than the handful of
tests shipped with the original problems, and the size of that difference was
measured directly on the task set used here: across the 63 tasks the benchmark's
own inputs number 193, against 6519 added by the augmentation. To ensure every task offers
structural headroom for the feedback loop to act on, tasks are filtered by an
objective, reproducible criterion applied to their reference solution: cyclomatic
complexity $\geq 3$ and $\geq 5$ lines of code, excluding tasks that require file
or network input/output. The filter is applied to the full pool of 378 tasks
distributed with the tool, is deterministic, and involves no sampling or manual
selection step: it yields exactly \textbf{63 tasks}. Because the criterion is
evaluated on the reference solutions, which are fixed properties of the benchmark,
the task set was settled before any code was generated and never revised
afterwards. These are stratified by the same
reference-solution complexity into a \emph{lower} group (cyclomatic complexity
3--4, $n=40$) and a \emph{higher} group (cyclomatic complexity $\geq 5$, $n=23$) so
that the effect of the loop can be analyzed by difficulty. The strata are therefore
a property of the benchmark tasks, fixed before any code was generated, and not of
the model's own baseline output. Correctness is decided by an
EvalPlus-style oracle: a solution passes only if, for every base and extended
input, its output matches the canonical reference solution (with floating-point
tolerance) and it raises no exception. The identical oracle is applied at every
iteration, which is how the study checks that code still works after
refinement: H3 is evaluated by comparing each task's baseline and final pass/fail
outcome under exactly the same inputs, and the McNemar
test~\citep{mcnemar1947} then operates on the
resulting discordant pairs (\cref{sec:results}). Each candidate is executed in an
isolated subprocess that is killed as a process group on timeout, with a second
in-process alarm inside the harness, so a non-terminating solution is recorded as a
failure rather than stalling the run.

Because H3 rests entirely on this oracle, the oracle itself was checked rather than
assumed. Submitting each task's own reference solution to it as though it were a
candidate, all 63 pass; submitting a deliberately corrupted variant of each
reference, all 63 fail. The augmented inputs also demonstrably discriminate: 40 of
the 63 baseline solutions satisfy the benchmark's original inputs, but 12 of those
are then caught by the augmentation, which is what reduces the baseline pass count
to the 28 reported in \cref{sec:results}. Almost a fifth of the task set would
therefore have been scored correct under the unaugmented benchmark.

\subsection{Loop Configurations}
\label{sec:configs}
Three configurations of the feedback loop are compared against the shared
baseline, all using the identical model, sampling parameters, task set, and
iteration cap. They differ only in the feedback prompt and in the policy that
decides which version is carried into the next iteration:
\begin{itemize}
  \item \textbf{Unguarded (comment-stripping).} The refactor prompt does not
    instruct the model to preserve existing comments. The label describes the
    observed outcome rather than the instruction: across the archived run the
    comment and docstring lines in the generated solutions fall from 268 at the
    baseline to 78 at the final iteration, against a rise to 324 under the
    comment-preserving prompt. Contrasting the two configurations tests whether
    comment handling affects the Maintainability Index.
  \item \textbf{Unguarded (comment-preserving).} Identical to the above except for
    one appended sentence instructing the model to keep any existing docstring and
    comments. Every other element of the prompt, including the requirement to preserve the
    exact function signature and behavior, is shared with the
    comment-stripping configuration, which is what makes the pair a controlled
    contrast. In both unguarded configurations the latest refactor is always
    carried forward, even if it fails the correctness check.
  \item \textbf{Correctness-guarded.} This configuration uses the
    comment-preserving prompt, so it differs from the second configuration in its
    carry-forward policy alone. The correctness check acts as a gate,
    mirroring a continuous-integration quality gate; the same principle as the
    fitness gate of Blyth et al.~\citep{blyth2025scam}, applied here to
    structural-metric refactoring rather than to linter
    violations. A refactor is accepted only
    if it still passes the oracle; otherwise it is rejected and the loop rolls
    back to the last accepted (still-correct) version, and the model is informed
    on the next attempt that its change altered behavior. That notice is not
    merely informative but structurally necessary: because decoding is greedy and
    the seed is fixed, rolling back and re-prompting with a byte-identical prompt
    would regenerate the same rejected refactor indefinitely. Adding the notice
    perturbs the prompt, which is what allows a retry to differ from the attempt it
    replaces. (The one exception is a
    task the baseline never solved: while the carried-forward version is still
    failing, any candidate is accepted, since correctness cannot regress below an already-failing baseline; this lets the loop repair such a task if it
    can.)
    By construction, final correctness can therefore never fall below the baseline.
\end{itemize}

A further comparison isolates the \emph{content} of the feedback rather than the
carry-forward policy. In all three configurations above, the model is shown its own
previous code together with the static-analysis findings. A \emph{code-blind}
variant instead withholds the previous code and supplies only the scalar findings
of the previous attempt---the Maintainability Index, the cyclomatic and cognitive
complexity totals, and the Pylint issue counts by category---asking the model to
generate a fresh, simpler solution to the same task. This tests whether the
improvement is driven by the model editing its own code or merely by being told
that its previous attempt was too complex. The code-blind variant is run in both an
unguarded and a correctness-guarded form and is reported as a secondary ablation
(\cref{sec:blind}).

\subsection{Frozen Protocol}
All generation parameters are fixed in a single version-controlled configuration
file and applied unchanged to every run: model \texttt{qwen2.5-coder:7b}
(Q4\_K\_M quantization), temperature~0, seed~42, a context window of 4096 tokens,
and an iteration cap of three. Because the model is served locally with greedy
decoding and a fixed seed, generations are not subject to provider drift, and the
generated code and per-iteration metrics are retained for every task and
iteration so that each reported value can be traced to the artifact that produced
it.

One limit on exact replication should be stated plainly: an Ollama tag identifies a
model by name rather than by content, so a future pull of \texttt{qwen2.5-coder:7b}
is not guaranteed to resolve to the same weights. The weights actually used are
nonetheless pinned by digest. The identifier \texttt{ollama list} reports for the
tag is \texttt{dae161e27b0e}, and the model layer it references is
\texttt{sha256:60e05f21\ldots} at 4.68\,GB; both are recorded in the configuration
file alongside the sampling parameters. These digests were read back from the local
model store after the experiment rather than logged during it, so the supporting
evidence that one set of weights served every run is that the manifest and the
weight blob both carry an unmodified timestamp preceding the earliest run; the tag
was pulled once and never re-pulled. A replication can therefore check that it
loaded the same weights by comparing these digests against its own, which the tag
alone would not allow.

\subsection{Artifacts and What Is Reproducible}
\label{sec:artifacts}
Two different claims are often both called reproducibility, and only the first of
them is exact here.

\emph{Recomputation from the archived artifacts is exact.} Every solution the model
produced is retained---1260 files across the five configurations, one per task and
iteration---together with the metric values recorded for it at run time. Every
number reported in this thesis is derived from those files by a deterministic
analysis script, and re-measuring all 1260 archived solutions reproduces all of the
recorded metric values exactly. A reader can therefore trace any reported figure to
the specific solution that produced it, and can confirm the statistics offline
without a model, a network connection, or an API key.

\emph{Re-running generation is approximate.} Single-shot generation is
deterministic: under greedy decoding with a fixed seed, the baseline iteration is
byte-identical across all five configurations. A multi-step trajectory is not
equally stable, because once any intermediate iteration differs the later ones can
diverge from it. A fresh run of the loop should therefore be expected to reproduce
the direction and rough magnitude of the effects reported here rather than the exact
values, and the figures quoted throughout are properly read as measurements of the
archived runs rather than as constants of the method. The consequences of this for
the study's conclusions are taken up in \cref{sec:limitations}.
